\chapter{Probabilistic STL (PrSTL)}
\label{ch:prstl}

Everything so far assumed the trace is the truth. It almost never is. A sensor reads with noise, a model has slack, and the same nominal run comes out a little different each time. A deterministic monitor answers \emph{does this exact trace satisfy the spec}. The question you usually have for real systems is \emph{how likely is the system to satisfy the spec}, and that is what PrSTL answers, with a probability and a confidence interval.

\section{The probabilistic operator}

PrSTL adds one operator. \code{P$\bowtie\alpha$ (phi)} asserts that \code{phi} holds with probability at least $\alpha$. The comparison $\bowtie$ can be any of $>, \ge, <, \le$, and the threshold $\alpha$ is a probability in $[0, 1]$ so \code{P>=1.5(x>0)} is a parse error. The body \code{phi} is any formula, which you can nest inside other temporal operators. The monitor estimates the probability by drawing many traces from the noise model and checking each one:

\begin{lstlisting}[language=Python]
Formula.parse("P>=0.9 (G (x > 0))")        # at least 0.9
Formula.parse("P>0.9 (G[0,10] (x > 5))")   # strictly above 0.9
Formula.parse("P<=0.1 (F (x < 60))")       # at most 0.1
Formula.parse("P<0.001 (F[0,5] (x < 0))")  # rare: under 1 in 1000
\end{lstlisting}

The body must be parenthesized so, \code{P>=0.9 x>0} is a parse error. And \code{P} is a reserved word, so you cannot have a signal named \code{P}.

\section{Fields of a probabilistic result}

To estimate the probability, \sentil{} needs to know how the readings vary, which you give it as a noise model per signal. Then \code{check} does the rest and hands back a result with every number you might want.

\begin{lstlisting}[language=Python]
import sentil
from sentil import Formula, LiftingRegistry, NoiseModel, SmcConfig

xs = [0.4 + 0.05 * i for i in range(20)]
trace = sentil.Trace(list(range(20)), {"x": xs})
lifting = LiftingRegistry()
lifting.register("x", NoiseModel.gaussian(0.0, 0.3))  # additive noise

phi = Formula.parse("P>=0.9 (G (x > 0))")
r = phi.check(trace, lifting, SmcConfig(samples=5000))

r.probability                 # 0.7332      fraction of draws that satisfied
r.satisfactions, r.samples    # 3666, 5000
r.interval.lower, r.interval.upper, r.interval.level # 0.7208, 0.7453, 0.95
r.holds                       # False       0.7332 does not meet >= 0.9
\end{lstlisting}

The seed defaults to $42$. The estimate is $0.73$: this rising signal clears zero comfortably on paper, but read through a sensor with a $0.3$ standard deviation, it slips under zero on more than a quarter of the draws, so it does not meet the $0.9$ bar.

\section{What check actually does}

\code{check} goes through five steps;

\begin{enumerate}[leftmargin=1.5em]
  \item \textbf{Lift.} For each of \code{samples} draws, copy the trace and replace every registered signal's readings with \code{reading + noise\_draw} (additive) or \code{reading * noise\_draw} (multiplicative). Unregistered signals pass through.
  \item \textbf{Score.} Evaluate the \emph{inner} formula, the body under \code{P}, on each lifted trace with the ordinary robustness engine. A draw satisfies when its robustness is at least zero.
  \item \textbf{Seed independently.} Draw $i$ is seeded from \code{seed + i}, so the count is identical.
  \item \textbf{Estimate.} \code{probability = satisfactions / samples}, here $3666/5000 = 0.7332$, and the confidence interval is built from the same counts by the Wilson method at the configured level. Wilson is the default because it is a good compromise between speed and accuracy, but \code{check\_conservative} uses the Clopper-Pearson exact interval for a more conservative estimate.
  \item \textbf{Decide.} \code{holds} compares the point estimate to the threshold in the operator's direction. $0.7332 \not\ge 0.9$, so \code{holds} is \code{False}.
\end{enumerate}

\section{The need for PrSTL}

If you run the same inner spec both ways on the same trace, you'll see the difference between the deterministic and probabilistic evaluations. Deterministically, this exact trace holds:

\begin{lstlisting}[language=Python]
Formula.parse("G (x > 0)").robustness(trace)   # 0.4  holds, set by x0
Formula.parse("P>=0.9 (G (x > 0))").check(trace, lifting).holds   # False
\end{lstlisting}

The deterministic monitor sees a positive margin and says yes. The $0.3$-standard-deviation sensor turns that thin $0.4$ margin into a coin that comes up violated more than a quarter of the time, so the probabilistic verdict is no. That gap, between \emph{this trace is fine} and \emph{the system is reliably fine}, is exactly what PrSTL is for.

\section{Confidence intervals}

\code{check\_conservative} swaps Wilson for the Clopper-Pearson exact interval when you must not overclaim, and \code{check\_distribution} returns the spread of robustness across the ensemble, which explains \emph{why} the probability is what it is.

\begin{lstlisting}[language=Python]
c = phi.check_conservative(trace, lifting, SmcConfig(samples=5000))
c.interval.lower, c.interval.upper           # 0.7207, 0.7454  (wider)

res, dist = phi.check_distribution(trace, lifting, SmcConfig(samples=5000))
dist.count, dist.mean, dist.std_dev, dist.min, dist.max
# 5000, 0.1037, 0.1881, -0.7163, 0.6583
\end{lstlisting}

The mean margin across the ensemble is barely positive, $0.10$, and the worst draw fell to $-0.72$. The property is holding by a hair. The distribution is the diagnosis behind the estimate.

\section{Why P is a verdict, not a margin}

Every other operator returns a robustness. \code{P} does not because a \code{P}-wrapped formula has no distance-to-satisfaction, because once the probability is estimated and compared to the threshold it is simply in or out. If you ask for its deterministic robustness, you get a SemanticError:

\begin{lstlisting}[language=Python]
Formula.parse("G (x > 0)").check(trace, lifting)   # SemanticError: needs P
Formula.parse("P>=0.9 (G (x > 0))").robustness(trace)
# SemanticError: the probabilistic operator P needs statistical evaluation;
#                deterministic robustness is undefined for it
\end{lstlisting}

The verdict for a \code{P} formula lives in \code{result.holds}.

\begin{pitfall}
\code{holds} compares the \emph{point} estimate to the threshold, which can flip near the line if you change the seed. If you want something more defensible, gate on the interval instead: accept \code{P>=0.9} only when \code{interval.lower >= 0.9}, and \code{P<0.001} only when \code{interval.upper < 0.001}. In the run above the whole interval sits below $0.9$, so both agree.
\end{pitfall}

\section{No noise, no probability}

One more thing, if you don't register a noise model, lifting is the identity, every draw is the same exact trace, and the estimate can only be $0.0$ or $1.0$:

\begin{lstlisting}[language=Python]
empty = LiftingRegistry()
empty.is_empty                                 # True
phi.check(trace, empty, SmcConfig(samples=5000)).probability   # 1.0
\end{lstlisting}

The five thousand samples did no work; the result is the deterministic verdict wearing a probability. Register a \code{NoiseModel} for every signal the inner formula touches, which is what the next chapter is about. Chapter~\ref{ch:smc} is then the statistics behind the interval, and the sequential and rare-event checks (\code{check\_sequential}, \code{check\_bayesian}, \code{check\_rare\_event}) for when a fixed sample budget is the wrong tool.